\begin{proof}

\begin{enumerate}[label={$\langle1\rangle$\arabic*.},leftmargin=*]
\item The evaluation, adjunction and fibre-equation comparisons
in (A.CD.1) use no Picard assumption and apply with
rational second-cohomology vanishing. They give a torsion
dualizing line $\omega_S$. The divisor Riemann--Roch
calculation gives $\chi(\OO_S)=0$. Thus, for
$e=h^0(S,\omega_S)\in\{0,1\}$, the same Serre-duality
calculation as in (A.CD.2) gives
$h^1(\OO_S)=1+e$; if $e=1$, then
$\omega_S\simeq\OO_S$.

\item Write $\eta:T\to S$ for normalization,
$\mathcal Q=\eta_*\OO_T/\OO_S$, and $\Gamma$ for
the full conductor image scheme. Equality of the
constant global functions on $S$ and $T$ gives the
actual injection $H^0(\mathcal Q)\hookrightarrow
H^1(\OO_S)$. Conductor duality and curve Serre duality
therefore give
\[
 h^1(\Gamma,\omega_S|_\Gamma)
       =h^0(S,\mathcal Q)\le1+e.                  \tag{A.J.31}\hypertarget{tag-0060}{}\label{cdp:eq:A.J.31}
\]
The restriction quotient
$\omega_S|_\Gamma\twoheadrightarrow\omega_S|_D$
has a kernel supported in dimension at most one, so
its $H^1$ map is surjective. For $e=1$ it follows that
$p_a(D)=h^1(\OO_D)\le2$. For $e=0$, put
$L=\omega_S|_D$. This torsion line has degree zero,
and \eqref{cdp:eq:A.J.31} gives $h^1(D,L)\le1$. Curve
Riemann--Roch gives
\[
 p_a(D)=1+h^1(D,L)-h^0(D,L)\le2,
\]
proving \eqref{cdp:eq:A.J.30} in both cases.

\item Let $\nu:\widehat D\to D$ be normalization. The
normalization sequence gives the exact numerical
comparison
\[
 p_a(D)=g(\widehat D)+
        \sum_{x\in D}\dim_{\C}
          (\nu_*\OO_{\widehat D}/\OO_D)_x.       \tag{A.J.32}\hypertarget{tag-0061}{}\label{cdp:eq:A.J.32}
\]
If $g(\widehat D)>0$, every local delta is consequently
at most one and their sum is at most one. We give the
analytic local rings at a point with delta one.
Let $A=\OO_{D,x}$ and let
$B=\bigoplus_{i=1}^r\C\{t_i\}$ be its finite
normalization. Here $\dim_{\C}(B/A)=1$. The residue map
of the local ring $A$ has the same value on every
normalization branch. The resulting quotient of
$\C^r$ by its diagonal has dimension $r-1$, so
$r\le2$.

\item For $r=2$, the subring
\[
 A_0=\{(b_1,b_2)\in\C\{t_1\}\oplus\C\{t_2\}:
                                  b_1(0)=b_2(0)\}
\]
contains $A$ and has colength one in $B$. Hence
$A=A_0$. The explicit homomorphism
$\C\{u,v\}/(uv)\to A_0$ sends
$u$ to $(t_1,0)$, $v$ to $(0,t_2)$, and a constant
to the corresponding equal pair. Expanding a
convergent series shows that it is an isomorphism.
This is the ordinary node germ.

\item For $r=1$, an element of $A$ with valuation one would
be a local coordinate on the normalization disc by
the holomorphic inverse function theorem. Convergent
composition in the analytic algebra $A$ would then
give $\C\{t\}\subset A$, contrary to its colength
one. Therefore every nonconstant element of $A$ has
valuation at least two. We have
$A\subset\C+t^2\C\{t\}$; both rings have
colength one in $\C\{t\}$, so equality holds.
The map
\[
 \C\{u,v\}/(v^2-u^3)\longrightarrow A,
                 \qquad (u,v)\longmapsto(t^2,t^3)
\]
is an isomorphism: every series with no linear term
is uniquely the sum of an even series in $t^2$ and
$t^3$ times a series in $t^2$. This is the ordinary
cusp germ. A delta-zero germ is its smooth
normalization disc.

\item We finally check the embedding into the specified
threefold. These analytic rings have embedding
dimension at most two. In local coordinates of the
smooth threefold, the tangent space of $D$ therefore
has codimension at least one. There is an element
$z$ of its ideal with $dz(x)\ne0$. The inverse
function theorem makes $z$ an ambient coordinate,
and the ideal of $D$ is the inverse image of a
height-one ideal in the regular analytic ring
$\C\{u,v\}=\OO_{Y,x}/(z)$. This ideal is principal:
a pure reduced curve in a smooth surface is a Cartier
divisor, or equivalently its reduced height-one ideal
is a product of the prime factors in this factorial
ring. Thus $I_{D,x}=(z,h)$, a regular sequence.
This proves the asserted lci embedding, including
at the node and cusp. Statement \ref{cdp:st:A.5.5} now applies
to the three specified generic transverse types
and proves the final smooth elliptic and constant-type
conclusion.

\end{enumerate}
\end{proof}
